\documentclass[10pt,a4paper]{amsart}
\usepackage{amsmath,amssymb,amsthm,graphicx}
% Retained minimal preamble for the Stacks Project chapter more-morphisms.tex (More on Morphisms).
% Verbatim excerpt of the chapter preamble: Xy-pic setup (native preamble lines 14-18) and the native
% theorem declarations with the mathematical macros (native preamble lines 154-217). Copied mechanically;
% nothing added, removed or redefined. The wrapper class and the amsmath/amssymb/amsthm/graphicx base
% packages are supplied by the assembler, so the class/externaldocument part of the native preamble is omitted.
% For commutative diagrams we use Xy-pic
\usepackage[all]{xy}

% We use 2cell for 2-commutative diagrams.
\xyoption{2cell}
% Theorem environments.
%
\theoremstyle{plain}
\newtheorem{theorem}[subsection]{Theorem}
\newtheorem{proposition}[subsection]{Proposition}
\newtheorem{lemma}[subsection]{Lemma}

\theoremstyle{definition}
\newtheorem{definition}[subsection]{Definition}
\newtheorem{example}[subsection]{Example}
\newtheorem{exercise}[subsection]{Exercise}
\newtheorem{situation}[subsection]{Situation}

\theoremstyle{remark}
\newtheorem{remark}[subsection]{Remark}
\newtheorem{remarks}[subsection]{Remarks}

\numberwithin{equation}{subsection}

% Macros
%
\def\lim{\mathop{\mathrm{lim}}\nolimits}
\def\colim{\mathop{\mathrm{colim}}\nolimits}
\def\Spec{\mathop{\mathrm{Spec}}}
\def\Hom{\mathop{\mathrm{Hom}}\nolimits}
\def\Ext{\mathop{\mathrm{Ext}}\nolimits}
\def\SheafHom{\mathop{\mathcal{H}\!\mathit{om}}\nolimits}
\def\SheafExt{\mathop{\mathcal{E}\!\mathit{xt}}\nolimits}
\def\Sch{\mathit{Sch}}
\def\Mor{\mathop{\mathrm{Mor}}\nolimits}
\def\Ob{\mathop{\mathrm{Ob}}\nolimits}
\def\Sh{\mathop{\mathit{Sh}}\nolimits}
\def\NL{\mathop{N\!L}\nolimits}
\def\CH{\mathop{\mathrm{CH}}\nolimits}
\def\proetale{{pro\text{-}\acute{e}tale}}
\def\etale{{\acute{e}tale}}
\def\QCoh{\mathit{QCoh}}
\def\Ker{\mathop{\mathrm{Ker}}}
\def\Im{\mathop{\mathrm{Im}}}
\def\Coker{\mathop{\mathrm{Coker}}}
\def\Coim{\mathop{\mathrm{Coim}}}

% Boxtimes
%
\DeclareMathSymbol{\boxtimes}{\mathbin}{AMSa}{"02}

%
% Macros for moduli stacks/spaces
%
\def\QCohstack{\mathcal{QC}\!\mathit{oh}}
\def\Cohstack{\mathcal{C}\!\mathit{oh}}
\def\Spacesstack{\mathcal{S}\!\mathit{paces}}
\def\Quotfunctor{\mathrm{Quot}}
\def\Hilbfunctor{\mathrm{Hilb}}
\def\Curvesstack{\mathcal{C}\!\mathit{urves}}
\def\Polarizedstack{\mathcal{P}\!\mathit{olarized}}
\def\Complexesstack{\mathcal{C}\!\mathit{omplexes}}
% \Pic is the operator that assigns to X its picard group, usage \Pic(X)
% \Picardstack_{X/B} denotes the Picard stack of X over B
% \Picardfunctor_{X/B} denotes the Picard functor of X over B
\def\Pic{\mathop{\mathrm{Pic}}\nolimits}
\def\Picardstack{\mathcal{P}\!\mathit{ic}}
\def\Picardfunctor{\mathrm{Pic}}
\def\Deformationcategory{\mathcal{D}\!\mathit{ef}}
\usepackage{cleveref}
\crefname{originalreference}{original reference}{original references}
\Crefname{originalreference}{Original reference}{Original references}
\makeatletter
\newlabel{morphisms-definition-family-ample-invertible-modules}{{0FXR}{source}{}{sourcepointer-1}{}}
\newlabel{morphisms-definition-family-ample-invertible-modules@cref}{{[originalreference][0][]0FXR}{[1][0][]source}{}{}{}}
\makeatother
\crefname{originalreference}{original reference}{original references}
\Crefname{originalreference}{Original reference}{Original references}
\makeatletter
\newlabel{properties-definition-ample}{{01PS}{source}{}{sourcepointer-2}{}}
\newlabel{properties-definition-ample@cref}{{[originalreference][0][]01PS}{[1][0][]source}{}{}{}}
\makeatother
\crefname{originalreference}{original reference}{original references}
\Crefname{originalreference}{Original reference}{Original references}
\makeatletter
\newlabel{properties-lemma-invert-s-sections}{{01PW}{source}{}{sourcepointer-3}{}}
\newlabel{properties-lemma-invert-s-sections@cref}{{[originalreference][0][]01PW}{[1][0][]source}{}{}{}}
\makeatother
\title{Effective Cartier divisors with ample restrictions produce an ample family}
\author{The Stacks Project Authors}
\date{stacks-project:more-morphisms.tex:local-snapshot-20261011; GNU FDL 1.2 or later}
\begin{document}
\maketitle
\noindent\textit{Editorial note.} Counted span is the author's own proof environment, source lines 23369-23389, immediately following the statement at 23360-23367.\par\medskip
\noindent\textit{Editorial note.} Three cross-chapter pointers (Morphisms, Properties) are recorded in assembly.external\_references.\par\medskip
\noindent\textit{Editorial note.} The author's step 'After replacing N by N + 1 we see that s' vanishes at every point of D\_i' is retained verbatim; it is the technical point of the proof.\par\medskip
\noindent\textit{Allowed standard prerequisites (editorial).}\begin{itemize}
\item Ampleness of an invertible module via affine non-vanishing loci. Properties, definition properties-definition-ample, cited at proof line 23375. Turns ampleness of L\_i restricted to U\_i into the existence of n and a section whose non-vanishing locus is affine.
\item Extending sections across an effective Cartier divisor. Properties, lemma properties-lemma-invert-s-sections, cited at proof lines 23376-23379. Extends the section from the complement of D\_i to all of X after twisting by a high multiple of O\_X(D\_i).
\item Definition of an ample family of invertible modules. Morphisms, definition morphisms-definition-family-ample-invertible-modules, cited at proof lines 23386-23388. Fixes the property that the constructed sections have to exhibit.
\end{itemize}\medskip
\ifdefined\hypertarget\hypertarget{sourcepointer-1}{}\fi
\noindent\textit{Source reference.} Original reference 0FXR: focused-seeds/geometry/S9-morph-get-ample-family/ref-tags.txt, source lines 5--5. Navigation pointer only: the counted proof cites this statement by \textbackslash{}ref as a standing input; the deduction the counted proof has to make is carried out inside the counted span, and no part of that deduction is delegated to this pointer.\par\medskip
\ifdefined\hypertarget\hypertarget{sourcepointer-2}{}\fi
\noindent\textit{Source reference.} Original reference 01PS: focused-seeds/geometry/S9-morph-get-ample-family/ref-tags.txt, source lines 6--6. Navigation pointer only: the counted proof cites this statement by \textbackslash{}ref as a standing input; the deduction the counted proof has to make is carried out inside the counted span, and no part of that deduction is delegated to this pointer.\par\medskip
\ifdefined\hypertarget\hypertarget{sourcepointer-3}{}\fi
\noindent\textit{Source reference.} Original reference 01PW: focused-seeds/geometry/S9-morph-get-ample-family/ref-tags.txt, source lines 7--7. Navigation pointer only: the counted proof cites this statement by \textbackslash{}ref as a standing input; the deduction the counted proof has to make is carried out inside the counted span, and no part of that deduction is delegated to this pointer.\par\medskip
\begin{lemma}
\label{lemma-get-ample-family}
Let $X$ be a scheme. Suppose given effective Cartier divisors
$D_1, \ldots, D_m$ on $X$ and invertible modules
$\mathcal{L}_1, \ldots, \mathcal{L}_m$ such that
$\bigcap D_i = \emptyset$ and $\mathcal{L}_i|_{X \setminus D_i}$
is ample. Then $X$ has an ample family of invertible modules.
\end{lemma}

\begin{proof}
Let $x \in X$. Choose an index $i \in \{1, \ldots, m\}$
such that $x \not \in D_i$. Set $U_i = X \setminus D_i$.
Since $\mathcal{L}_i|_{U_i}$ we can find an $n \geq 1$ and
a section $s \in \Gamma(U_i, \mathcal{L}_i^{\otimes n})$
such that the locus $(U_i)_s$ where $s$ doesn't vanish is affine
(Properties, Definition \ref{properties-definition-ample}).
Since $U_i$ is the locus where the canonical section
$1 \in \mathcal{O}_X(D_i)$ doesn't vanish, we see from
Properties, Lemma \ref{properties-lemma-invert-s-sections}
there exists an $N \geq 0$ such that $s$ extends to a section
$$
s' \in \Gamma(X, \mathcal{L}_i^{\otimes n} \otimes_{\mathcal{O}_X}
\mathcal{O}_X(N D_i))
$$
After replacing $N$ by $N + 1$ we see that $s'$ vanishes at every
point of $D_i$ and hence that $X_{s'} = (U_i)_s$ is affine.
This proves that $X$ has an ample family of invertible modules, see
Morphisms, Definition
\ref{morphisms-definition-family-ample-invertible-modules}.
\end{proof}

\end{document}
